\documentclass[11pt,a4paper]{article}
\usepackage[utf8]{inputenc}
\usepackage[T1]{fontenc}
\usepackage{amsmath,amssymb,amsthm}
\usepackage{listings}
\usepackage{geometry}
\usepackage{hyperref}
\geometry{margin=2.5cm}

\lstset{
  language=C++,
  basicstyle=\ttfamily\small,
  keywordstyle=\bfseries,
  breaklines=true,
  numbers=left,
  numberstyle=\tiny,
  frame=single,
  tabsize=2
}

\title{IOI 2021 -- Candies}
\author{Solution}
\date{}

\begin{document}
\maketitle

\section{Problem Summary}
There are $n$ boxes, each with capacity $c[i]$. There are $q$ operations, where operation $j$ adds $v[j]$ candies to all boxes in range $[l[j], r[j]]$. If $v[j] > 0$, candies are added (capped at capacity). If $v[j] < 0$, candies are removed (capped at 0). After all operations, find the final number of candies in each box.

\section{Solution Approach}

\subsection{Key Insight: Process in Reverse}
For each box $i$, consider the sequence of operations that affect it (in order). The final value depends on the \emph{last capping event}: the last time the box either completely filled or completely emptied. After that event, the remaining operations just add/subtract without capping.

\subsection{Algorithm}
\begin{enumerate}
  \item Process operations in reverse. Maintain the prefix sums of $v$ values.
  \item For each box $i$ with capacity $c[i]$, find the most recent operation $j$ such that the running sum from operation $j$ onward exceeds $c[i]$ (cap at capacity) or goes below $0$ (cap at zero).
  \item After finding this capping point, the final value is either $c[i] - (\text{excess above } c)$ or $0 + (\text{amount above } 0)$.
\end{enumerate}

\subsection{Segment Tree Approach}
Use a segment tree over the operations (reversed). For each box $i$:
\begin{itemize}
  \item Find the latest operation $j$ where the prefix sum from $j$ to $q$ has maximum $> c[i]$ or minimum $< 0$.
  \item The segment tree stores prefix sums and supports range max/min queries.
  \item Binary search on the segment tree to find the critical operation.
\end{itemize}

Concretely:
\begin{enumerate}
  \item Compute suffix sums $S[j] = \sum_{k=j}^{q-1} v[k]$ for operations affecting box $i$.
  \item Find the latest $j$ such that $\max(S[j..q]) - \min(S[j..q]) > c[i]$.
  \item If no such $j$ exists, the box never caps: final value = $\max(0, \min(c[i], S[0]))$.
  \item Otherwise, at operation $j$: if the extreme is a maximum, the box capped at full, and remaining value is $c[i] - (\max_{k \ge j} S[k] - S[\text{end}])$. If the extreme is a minimum, it capped at empty.
\end{enumerate}

For efficiency, we avoid processing each box independently. Instead, use a sweep:
\begin{itemize}
  \item Sort boxes by capacity and process using the segment tree.
  \item Or, for each box, binary-search on the segment tree of all operations (filtering to only those affecting box $i$ is handled by building the segment tree with events at appropriate positions).
\end{itemize}

The efficient approach processes each box in $O(\log q)$ time using binary search on a global segment tree of prefix sums.

\section{C++ Solution}

\begin{lstlisting}
#include <bits/stdc++.h>
using namespace std;
typedef long long ll;

struct SegTree {
    int n;
    vector<ll> mx, mn; // max and min of prefix sums in range

    SegTree(int n) : n(n), mx(4 * n, LLONG_MIN), mn(4 * n, LLONG_MAX) {}

    void update(int pos, ll val, int node, int lo, int hi) {
        if (lo == hi) {
            mx[node] = mn[node] = val;
            return;
        }
        int mid = (lo + hi) / 2;
        if (pos <= mid) update(pos, val, 2*node, lo, mid);
        else update(pos, val, 2*node+1, mid+1, hi);
        mx[node] = max(mx[2*node], mx[2*node+1]);
        mn[node] = min(mn[2*node], mn[2*node+1]);
    }

    void update(int pos, ll val) { update(pos, val, 1, 0, n-1); }

    // Find the latest position in [lo, hi] where max - min of suffix > cap
    // Returns the position or -1
    // Actually, we need a different query: find latest j such that
    // max(S[j..q]) - min(S[j..q]) > c

    // Query: max in range [l, r]
    ll query_max(int l, int r, int node, int lo, int hi) {
        if (l > hi || r < lo) return LLONG_MIN;
        if (l <= lo && hi <= r) return mx[node];
        int mid = (lo + hi) / 2;
        return max(query_max(l, r, 2*node, lo, mid),
                   query_max(l, r, 2*node+1, mid+1, hi));
    }
    ll query_max(int l, int r) { return query_max(l, r, 1, 0, n-1); }

    ll query_min(int l, int r, int node, int lo, int hi) {
        if (l > hi || r < lo) return LLONG_MAX;
        if (l <= lo && hi <= r) return mn[node];
        int mid = (lo + hi) / 2;
        return min(query_min(l, r, 2*node, lo, mid),
                   query_min(l, r, 2*node+1, mid+1, hi));
    }
    ll query_min(int l, int r) { return query_min(l, r, 1, 0, n-1); }
};

vector<int> distribute_candies(vector<int> c, vector<int> l,
                                 vector<int> r, vector<int> v) {
    int n = c.size(), q = l.size();
    vector<int> ans(n);

    // For each operation j, it affects boxes l[j]..r[j].
    // We process boxes using sweep: for each box i, collect operations affecting it.
    // Use event-based approach: at l[j], operation j starts; at r[j]+1, it ends.

    // Build prefix sums of v values over operations applied to box i.
    // Instead, use a segment tree over operations and process each box.

    // Events: operations sorted by l and r
    vector<vector<int>> starts(n), ends(n);
    for (int j = 0; j < q; j++) {
        starts[l[j]].push_back(j);
        if (r[j] + 1 < n) ends[r[j] + 1].push_back(j);
    }

    // Segment tree over operations [0..q] where position j holds prefix sum S[j]
    // S[j] = sum of v[0..j-1] for operations active at current box
    // We maintain S[j] = cumulative sum up to operation j.

    // Actually, let's define the prefix sum differently.
    // For box i, let the operations affecting it be o_1, o_2, ..., o_t (in order).
    // Prefix sum P[0] = 0, P[k] = P[k-1] + v[o_k].
    // The final value depends on the latest k such that max(P[k..t]) - min(P[k..t]) >= c[i].

    // Global approach: maintain a segment tree where position j has value S[j] if
    // operation j is active (affects current box i), and -inf otherwise.
    // As we sweep i from 0 to n-1, we add/remove operations.

    // S[j] = running prefix sum of active operations up to j.
    // This is complex to maintain dynamically.

    // Alternative cleaner approach: segment tree with all q+1 positions.
    // S[0] = 0. For j = 1..q: S[j] = S[j-1] + (v[j-1] if op j-1 is active else 0).

    // When sweeping boxes, adding/removing an operation j changes S[j+1], S[j+2], ..., S[q]
    // by +/- v[j]. This is a range update on the prefix sums.

    // Segment tree with lazy propagation supporting:
    //   - Range add on [j+1, q]
    //   - Query max and min on suffix [j, q]
    //   - Binary search for latest j where max[j..q] - min[j..q] > c[i]

    // This is feasible but complex. Let me implement it.

    // Segment tree with lazy add, range max/min query
    struct LazySegTree {
        int n;
        vector<ll> mx, mn, lazy;

        LazySegTree(int n) : n(n), mx(4*(n+1), 0), mn(4*(n+1), 0), lazy(4*(n+1), 0) {
            // Initially all values are 0
        }

        void push_down(int node) {
            if (lazy[node] != 0) {
                for (int c : {2*node, 2*node+1}) {
                    mx[c] += lazy[node];
                    mn[c] += lazy[node];
                    lazy[c] += lazy[node];
                }
                lazy[node] = 0;
            }
        }

        void range_add(int l, int r, ll val, int node, int lo, int hi) {
            if (l > hi || r < lo) return;
            if (l <= lo && hi <= r) {
                mx[node] += val;
                mn[node] += val;
                lazy[node] += val;
                return;
            }
            push_down(node);
            int mid = (lo + hi) / 2;
            range_add(l, r, val, 2*node, lo, mid);
            range_add(l, r, val, 2*node+1, mid+1, hi);
            mx[node] = max(mx[2*node], mx[2*node+1]);
            mn[node] = min(mn[2*node], mn[2*node+1]);
        }
        void range_add(int l, int r, ll val) { range_add(l, r, val, 1, 0, n); }

        ll qmax(int l, int r, int node, int lo, int hi) {
            if (l > hi || r < lo) return LLONG_MIN;
            if (l <= lo && hi <= r) return mx[node];
            push_down(node);
            int mid = (lo + hi) / 2;
            return max(qmax(l, r, 2*node, lo, mid), qmax(l, r, 2*node+1, mid+1, hi));
        }
        ll qmax(int l, int r) { return qmax(l, r, 1, 0, n); }

        ll qmin(int l, int r, int node, int lo, int hi) {
            if (l > hi || r < lo) return LLONG_MAX;
            if (l <= lo && hi <= r) return mn[node];
            push_down(node);
            int mid = (lo + hi) / 2;
            return min(qmin(l, r, 2*node, lo, mid), qmin(l, r, 2*node+1, mid+1, hi));
        }
        ll qmin(int l, int r) { return qmin(l, r, 1, 0, n); }

        // Find the latest position j in [0, q] such that
        // max(S[j..q]) - min(S[j..q]) > cap
        // Binary search: start from q, expand left
        int find_latest(ll cap) {
            // Binary search on j from q down to 0
            ll cur_max = LLONG_MIN, cur_min = LLONG_MAX;
            // Start from j = q (rightmost), check j = q, q-1, ...
            // When max - min > cap, return j.
            // Use segment tree descent for efficiency.

            // Simple approach: binary search on j
            int lo = 0, hi = n, result = -1;
            while (lo <= hi) {
                int mid = (lo + hi) / 2;
                ll mx_val = qmax(mid, n);
                ll mn_val = qmin(mid, n);
                if (mx_val - mn_val > cap) {
                    result = mid;
                    lo = mid + 1; // try later (larger j)
                } else {
                    hi = mid - 1;
                }
            }
            return result;
        }

        ll point_query(int pos, int node, int lo, int hi) {
            if (lo == hi) return mx[node];
            push_down(node);
            int mid = (lo + hi) / 2;
            if (pos <= mid) return point_query(pos, 2*node, lo, mid);
            else return point_query(pos, 2*node+1, mid+1, hi);
        }
        ll point_query(int pos) { return point_query(pos, 1, 0, n); }
    };

    LazySegTree seg(q);
    // seg represents S[0..q] where S[0] = 0 initially.
    // S[j] = prefix sum of v values for operations 0..j-1 that affect current box.

    for (int i = 0; i < n; i++) {
        // Add operations starting at i
        for (int j : starts[i]) {
            seg.range_add(j + 1, q, v[j]);
        }
        // Remove operations ending before i
        for (int j : ends[i]) {
            seg.range_add(j + 1, q, -v[j]);
        }

        ll cap = c[i];
        int j = seg.find_latest(cap);

        ll S_end = seg.point_query(q); // S[q] = total sum

        if (j == -1) {
            // No capping event
            ans[i] = (int)max(0LL, min((ll)cap, S_end));
        } else {
            // Capping happened at operation j
            // Check if it was a max or min event
            ll mx_from_j = seg.qmax(j, q);
            ll mn_from_j = seg.qmin(j, q);

            // The latest j where the range exceeds cap.
            // At j, either the max caused overflow (cap at full) or min caused underflow (cap at 0)
            // Need to check which extreme at j+1 triggers:
            // If max(S[j..q]) - S[q] part... let me think.
            // After operation j, if max was reached: box was full, final = cap - (max - S_end)
            // If min was reached: box was empty, final = 0 + (S_end - min)

            // The capping event at j:
            // S[j] is the prefix sum just before operation j.
            // From j onward, the range of S values is [mn_from_j, mx_from_j].
            // If the max of S[j..q] - mn of S[j..q] > cap, one of the extremes causes capping.

            // After the last capping: if max was hit last (later than min), final = cap - (mx - S_end)
            // If min was hit last, final = S_end - mn.
            // But we need to know which extreme was reached later.

            // Find within [j, q]: position of max and position of min.
            // The one later determines the capping direction.

            // For simplicity: check S[j] value.
            // If S[j] is closer to the max: the step from j-1 to j pushed upward -> max event
            // Actually, more precisely: we need to check from j+1:
            // max(S[j+1..q]) - min(S[j+1..q]) <= cap (by definition of j being latest)
            // So the capping is at operation j exactly.

            // S[j] is either the new max or new min.
            ll S_j = seg.point_query(j);
            ll mx_after = seg.qmax(j+1, q);
            ll mn_after = seg.qmin(j+1, q);

            // Adding S[j] to the range [mn_after, mx_after] creates overflow.
            if (S_j >= mx_after) {
                // S[j] is the max, causing cap at full
                // After cap: box = cap, then operations j+1..q bring it to:
                // final = cap + (S_end - S_j)... but capping at cap means
                // the box is at cap after operation j.
                // Then subsequent operations: box = cap + (S_end - S_j)
                // but clamped... since j is the LATEST capping, no more clamping after.
                ans[i] = (int)(cap - (S_j - S_end));
            } else {
                // S[j] is the min, causing cap at 0
                ans[i] = (int)(S_end - S_j);
            }
            ans[i] = max(0, min((int)cap, ans[i]));
        }
    }

    return ans;
}

int main() {
    int n, q;
    scanf("%d %d", &n, &q);
    vector<int> c(n);
    for (int i = 0; i < n; i++) scanf("%d", &c[i]);
    vector<int> l(q), r(q), v(q);
    for (int j = 0; j < q; j++) scanf("%d %d %d", &l[j], &r[j], &v[j]);
    auto ans = distribute_candies(c, l, r, v);
    for (int i = 0; i < n; i++)
        printf("%d%c", ans[i], " \n"[i == n-1]);
    return 0;
}
\end{lstlisting}

\section{Complexity Analysis}
\begin{itemize}
  \item \textbf{Time complexity:} $O((n + q) \log q)$ -- each box requires $O(\log q)$ for the binary search on the segment tree, and maintaining the segment tree as we sweep takes $O(\log q)$ per operation add/remove.
  \item \textbf{Space complexity:} $O(n + q)$.
\end{itemize}

\end{document}
